\documentclass[12pt,a4paper]{ctexart}

% 页面设置
\usepackage[
    left=2.5cm,
    right=2.5cm,
    top=2.3cm,
    bottom=2.3cm
]{geometry}

% 常用宏包
\usepackage{xcolor}
\usepackage{tcolorbox}
\usepackage{titlesec}
\usepackage{fancyhdr}
\usepackage{enumitem}
\usepackage{tabularx}
\usepackage{amsmath, amssymb}
\usepackage{hyperref}

\tcbuselibrary{skins, breakable}

% ======================
% 颜色主题：清爽学习笔记风
% ======================

% 主色：深蓝绿
\definecolor{MainColor}{HTML}{2A6F73}

% 辅助色：青绿色
\definecolor{SecondColor}{HTML}{3CAEA3}

% 浅背景色
\definecolor{LightMain}{HTML}{EAF7F6}

% 标题深色
\definecolor{TitleDark}{HTML}{1F3A40}

% 文字灰
\definecolor{TextGray}{HTML}{555555}

% 边框灰
\definecolor{BorderGray}{HTML}{CBD5D6}

% 提醒色：柔和橙
\definecolor{WarnColor}{HTML}{F2B35D}

% 易错点红色
\definecolor{ErrorColor}{HTML}{C94C4C}

% 定义框紫色
\definecolor{DefineColor}{HTML}{6C63B7}

% ======================
% 正文设置
% ======================

\linespread{1.25}
\setlength{\parindent}{2em}
\setlength{\parskip}{0.25em}

\hypersetup{
  colorlinks=true,
  linkcolor=MainColor,
  urlcolor=MainColor
}

% ======================
% 页眉页脚
% ======================

\pagestyle{fancy}
\fancyhf{}
\fancyfoot[L]{\textcolor{MainColor}{机器学习笔记}}
\fancyfoot[R]{\textcolor{MainColor}{\thepage}}
\renewcommand{\headrulewidth}{0pt}
\renewcommand{\footrulewidth}{0pt}

% ======================
% 顶部小标题栏
% ======================

\newcommand{\topbar}[2]{
  \noindent
  {\small\textcolor{MainColor}{#1}}
  \hfill
  {\small\bfseries\textcolor{TitleDark}{#2}}
  \par\vspace{0.35em}
  {\color{BorderGray}\hrule}
  \vspace{1.2em}
}

% ======================
% 标题卡片
% ======================

\newcommand{\notetitlecard}[6]{
\begin{tcolorbox}[
  enhanced,
  colback=white,
  colframe=BorderGray,
  boxrule=0.6pt,
  arc=2mm,
  left=12pt,
  right=12pt,
  top=10pt,
  bottom=10pt,
  borderline west={4pt}{0pt}{SecondColor}
]
\begin{tabularx}{\linewidth}{X r}
{\small\textcolor{TextGray}{#1}} & {\small\textcolor{TextGray}{#2}} \\[0.9em]
\multicolumn{2}{l}{
  {\Huge\bfseries\textcolor{MainColor}{#3}}
} \\[0.45em]
\multicolumn{2}{l}{
  {\large\textcolor{SecondColor!90!black}{#4}}
} \\[0.45em]
\multicolumn{2}{l}{
  {\small\textcolor{TextGray}{课程：#5 \qquad 作者：#6}}
}
\end{tabularx}
\end{tcolorbox}
\vspace{1em}
}

% ======================
% 摘要框
% ======================

\newtcolorbox{summarybox}[1][本节摘要]{
  enhanced,
  breakable,
  colback=LightMain,
  colframe=SecondColor,
  colbacktitle=SecondColor,
  coltitle=white,
  title={#1},
  fonttitle=\bfseries,
  boxrule=0.6pt,
  arc=2mm,
  left=12pt,
  right=12pt,
  top=10pt,
  bottom=10pt
}

% ======================
% 定义框
% ======================

\newtcolorbox{definitionbox}[1][定义]{
  enhanced,
  breakable,
  colback=DefineColor!6,
  colframe=DefineColor,
  colbacktitle=DefineColor,
  coltitle=white,
  title={#1},
  fonttitle=\bfseries,
  boxrule=0.6pt,
  arc=2mm,
  left=12pt,
  right=12pt,
  top=10pt,
  bottom=10pt
}

% ======================
% 重点框
% ======================

\newtcolorbox{keybox}[1][重点]{
  enhanced,
  breakable,
  colback=MainColor!5,
  colframe=MainColor,
  colbacktitle=MainColor,
  coltitle=white,
  title={#1},
  fonttitle=\bfseries,
  boxrule=0.6pt,
  arc=2mm,
  left=12pt,
  right=12pt,
  top=10pt,
  bottom=10pt
}

% ======================
% 例题框
% ======================

\newtcolorbox{examplebox}[1][例题]{
  enhanced,
  breakable,
  colback=white,
  colframe=SecondColor,
  colbacktitle=SecondColor!85!black,
  coltitle=white,
  title={#1},
  fonttitle=\bfseries,
  boxrule=0.6pt,
  arc=2mm,
  left=12pt,
  right=12pt,
  top=10pt,
  bottom=10pt
}

% ======================
% 易错点框
% ======================

\newtcolorbox{mistakebox}[1][易错点]{
  enhanced,
  breakable,
  colback=ErrorColor!5,
  colframe=ErrorColor,
  colbacktitle=ErrorColor,
  coltitle=white,
  title={#1},
  fonttitle=\bfseries,
  boxrule=0.6pt,
  arc=2mm,
  left=12pt,
  right=12pt,
  top=10pt,
  bottom=10pt
}

% ======================
% 提醒框
% ======================

\newtcolorbox{tipbox}[1][提示]{
  enhanced,
  breakable,
  colback=WarnColor!10,
  colframe=WarnColor,
  colbacktitle=WarnColor,
  coltitle=white,
  title={#1},
  fonttitle=\bfseries,
  boxrule=0.6pt,
  arc=2mm,
  left=12pt,
  right=12pt,
  top=10pt,
  bottom=10pt
}

% ======================
% 章节标题样式
% ======================

\titleformat{\section}
  {\Large\bfseries\color{MainColor}}
  {\thesection}
  {0.8em}
  {}

\titleformat{\subsection}
  {\large\bfseries\color{TitleDark}}
  {\thesubsection}
  {0.8em}
  {}

\titleformat{\subsubsection}
  {\normalsize\bfseries\color{SecondColor!80!black}}
  {\thesubsubsection}
  {0.8em}
  {}

% ======================
% 列表样式
% ======================

\setlist[itemize]{
  leftmargin=2.2em,
  itemsep=0.3em,
  topsep=0.4em
}

\setlist[enumerate]{
  leftmargin=2.2em,
  itemsep=0.3em,
  topsep=0.4em
}

% ======================
% 自定义强调命令
% ======================

\newcommand{\important}[1]{\textbf{\textcolor{MainColor}{#1}}}
\newcommand{\warning}[1]{\textbf{\textcolor{ErrorColor}{#1}}}
\newcommand{\concept}[1]{\textbf{\textcolor{DefineColor}{#1}}}

% ======================
% 正文开始
% ======================

\begin{document}

\topbar{吴恩达机器学习}{学习笔记}

\notetitlecard
  {Study Notes Series}
  {2026 Spring}
  {吴恩达机器学习笔记(3)}
  {}
  {机器学习}
  {C++与草稿纸}

\begin{summarybox}
本节围绕分类问题与过拟合问题，介绍逻辑回归从模型表示、代价函数到参数优化的完整流程，并说明如何使用正则化提升模型对新数据的泛化能力。
\begin{itemize}[leftmargin=2em, itemsep=2pt, topsep=4pt]
  \item \textbf{逻辑回归：}使用 Sigmoid 函数把线性组合映射为 $0$ 到 $1$ 之间的概率；
  \item \textbf{判定边界：}由参数决定输入空间中不同预测类别的分界，可通过多项式特征形成非线性边界；
  \item \textbf{模型训练：}使用二元交叉熵代价函数，并通过梯度下降或高级优化算法求解参数；
  \item \textbf{多类别分类：}训练多个“一对多”二元分类器，选择预测概率最大的类别；
  \item \textbf{正则化：}在代价函数中惩罚过大的参数，缓解过拟合，同时通过参数 $\lambda$ 控制拟合与简化模型之间的平衡。
\end{itemize}
\end{summarybox}

\section{分类与逻辑回归}

\subsection{分类问题}

\begin{definitionbox}[二元分类]
分类问题的目标是预测离散类别，而不是连续数值。对于二元分类，通常把两个类别分别记为
\[
  y\in\{0,1\},
\]
其中 $0$ 表示负类（negative class），$1$ 表示正类（positive class）。例如，判断邮件是否为垃圾邮件、交易是否存在欺诈、肿瘤是良性还是恶性，都可以写成二元分类问题。
\end{definitionbox}

\begin{mistakebox}[为什么不直接使用线性回归分类]
线性回归的输出没有范围限制，可能小于 $0$ 或大于 $1$。即使人为规定“预测值不小于 $0.5$ 时判为正类”，新加入的极端样本也可能明显改变拟合直线和分类结果。因此，线性回归的输出既不适合解释为概率，也通常不是分类任务的合理模型。

逻辑回归虽然名称中含有“回归”，但它本质上是一个\important{分类算法}。
\end{mistakebox}

\subsection{Sigmoid 函数与假设表示}

\begin{definitionbox}[Sigmoid 函数]
逻辑回归使用 Sigmoid 函数，也称 Logistic 函数：
\[
  g(z)=\frac{1}{1+e^{-z}}.
\]
它把任意实数映射到区间 $(0,1)$，并具有以下性质：
\[
  g(0)=0.5,\qquad
  z>0\Rightarrow g(z)>0.5,\qquad
  z<0\Rightarrow g(z)<0.5.
\]
当 $z$ 趋向正无穷时，$g(z)$ 趋近于 $1$；当 $z$ 趋向负无穷时，$g(z)$ 趋近于 $0$。
\end{definitionbox}

\begin{keybox}[逻辑回归模型]
令 $x_0=1$，特征向量和参数向量分别为
\[
  \mathbf{x}=[x_0,x_1,\ldots,x_n]^T,
  \qquad
  \boldsymbol{\theta}=[\theta_0,\theta_1,\ldots,\theta_n]^T.
\]
逻辑回归模型为
\[
  h_{\boldsymbol{\theta}}(\mathbf{x})
  =g\!\left(\boldsymbol{\theta}^{T}\mathbf{x}\right)
  =\frac{1}{1+e^{-\boldsymbol{\theta}^{T}\mathbf{x}}}.
\]
模型输出可以解释为在给定输入和参数时属于正类的估计概率：
\[
  h_{\boldsymbol{\theta}}(\mathbf{x})
  =P(y=1\mid\mathbf{x};\boldsymbol{\theta}).
\]
因此
\[
  P(y=0\mid\mathbf{x};\boldsymbol{\theta})
  =1-h_{\boldsymbol{\theta}}(\mathbf{x}).
\]
例如，若 $h_{\boldsymbol{\theta}}(\mathbf{x})=0.7$，则模型认为该样本属于正类的概率为 $70\%$。
\end{keybox}

\begin{tipbox}[概率与最终类别]
模型输出的是概率，最终类别还需要通过阈值获得。默认阈值通常取 $0.5$：
\[
  \hat y=
  \begin{cases}
    1, & h_{\boldsymbol{\theta}}(\mathbf{x})\geq 0.5,\\
    0, & h_{\boldsymbol{\theta}}(\mathbf{x})<0.5.
  \end{cases}
\]
实际应用中可以根据漏判与误判的代价调整阈值；改变阈值会改变分类结果，但不会重新训练模型参数。
\end{tipbox}

\subsection{判定边界}

\begin{definitionbox}[判定边界]
当阈值为 $0.5$ 时，Sigmoid 函数的单调性给出
\[
  h_{\boldsymbol{\theta}}(\mathbf{x})\geq 0.5
  \Longleftrightarrow
  \boldsymbol{\theta}^{T}\mathbf{x}\geq 0.
\]
所以方程
\[
  \boldsymbol{\theta}^{T}\mathbf{x}=0
\]
描述的集合就是\concept{判定边界（Decision Boundary）}。判定边界由模型参数和特征表示共同决定，而不是训练样本本身直接构成。
\end{definitionbox}

\begin{examplebox}[线性判定边界]
若只有两个特征，并且
\[
  \boldsymbol{\theta}=
  \begin{bmatrix}-3&1&1\end{bmatrix}^{T},
\]
则
\[
  \boldsymbol{\theta}^{T}\mathbf{x}=-3+x_1+x_2.
\]
判定边界为直线 $x_1+x_2=3$。直线一侧满足 $x_1+x_2\geq3$，预测为 $1$；另一侧预测为 $0$。
\end{examplebox}

\begin{examplebox}[非线性判定边界]
逻辑回归对参数仍然是线性的，但可以使用多项式特征表示弯曲边界。例如
\[
  h_{\boldsymbol{\theta}}(\mathbf{x})
  =g\!\left(-1+x_1^2+x_2^2\right)
\]
对应的判定边界为
\[
  x_1^2+x_2^2=1,
\]
即以原点为圆心、半径为 $1$ 的圆。加入更多高次项或交叉项可以得到更复杂的边界，但模型越复杂，过拟合风险通常越高。
\end{examplebox}

\section{逻辑回归的代价函数与训练}

\subsection{二元交叉熵代价函数}

\begin{mistakebox}[不能直接照搬平方误差]
将 Sigmoid 模型代入线性回归的平方误差代价函数，会得到不利于梯度下降求解的非凸目标。逻辑回归改用对数损失，使代价函数成为凸函数，从而避免多个局部最小值带来的优化困难。
\end{mistakebox}

\begin{definitionbox}[单个样本的对数损失]
令 $\hat y=h_{\boldsymbol{\theta}}(\mathbf{x})$，则一个训练样本的损失为
\[
  L(\hat y,y)=
  \begin{cases}
    -\log(\hat y), & y=1,\\
    -\log(1-\hat y), & y=0.
  \end{cases}
\]
两种情况可以合并为
\[
  L(\hat y,y)
  =-y\log(\hat y)-(1-y)\log(1-\hat y).
\]
当预测正确且置信度高时，损失接近 $0$；当模型以很高置信度给出错误预测时，损失会迅速增大。
\end{definitionbox}

\begin{keybox}[训练集上的代价函数]
对 $m$ 个训练样本取平均，逻辑回归的代价函数为
\[
  J(\boldsymbol{\theta})
  =-\frac{1}{m}\sum_{i=1}^{m}
  \left[
    y^{(i)}\log h_{\boldsymbol{\theta}}\!\left(\mathbf{x}^{(i)}\right)
    +(1-y^{(i)})\log\!\left(1-h_{\boldsymbol{\theta}}\!\left(\mathbf{x}^{(i)}\right)\right)
  \right].
\]
训练目标是求
\[
  \underset{\boldsymbol{\theta}}{\operatorname{minimize}}
  \;J(\boldsymbol{\theta}).
\]
这个损失也称为\concept{二元交叉熵（Binary Cross-Entropy）}。
\end{keybox}

\subsection{梯度下降}

\begin{definitionbox}[逻辑回归的梯度]
对每个参数 $\theta_j$，代价函数的偏导数为
\[
  \frac{\partial J(\boldsymbol{\theta})}{\partial\theta_j}
  =\frac{1}{m}\sum_{i=1}^{m}
  \left(
    h_{\boldsymbol{\theta}}\!\left(\mathbf{x}^{(i)}\right)-y^{(i)}
  \right)x_j^{(i)},
  \qquad j=0,1,\ldots,n.
\]
梯度下降同时更新所有参数：
\[
  \theta_j\leftarrow\theta_j
  -\alpha\frac{1}{m}\sum_{i=1}^{m}
  \left(
    h_{\boldsymbol{\theta}}\!\left(\mathbf{x}^{(i)}\right)-y^{(i)}
  \right)x_j^{(i)}.
\]
它在形式上与线性回归的更新式相似，但两者的假设函数不同：线性回归使用 $\boldsymbol{\theta}^T\mathbf{x}$，逻辑回归使用 $g(\!\boldsymbol{\theta}^T\mathbf{x})$。
\end{definitionbox}

\begin{tipbox}[向量化与特征缩放]
令特征矩阵 $X\in\mathbb{R}^{m\times(n+1)}$、标签向量 $\mathbf{y}\in\mathbb{R}^{m}$，全部样本的预测向量为
\[
  \hat{\mathbf{y}}=g(X\boldsymbol{\theta}),
\]
其中 $g$ 对向量中的每个元素分别计算。于是
\[
  \nabla J(\boldsymbol{\theta})
  =\frac{1}{m}X^T(\hat{\mathbf{y}}-\mathbf{y}).
\]
逻辑回归同样受特征尺度影响。训练前进行特征缩放，通常能使梯度下降更快收敛。
\end{tipbox}

\begin{mistakebox}[实现交叉熵时的数值问题]
由于 $\log(0)$ 没有有限值，直接计算概率形式的交叉熵可能在浮点数下产生无穷大或非数。实践中应使用数值稳定的库函数，或在计算对数前把概率限制在 $(\varepsilon,1-\varepsilon)$ 内。该处理只用于避免数值溢出，不应改变损失函数的含义。
\end{mistakebox}

\subsection{高级优化算法}

\begin{tipbox}[梯度下降以外的选择]
共轭梯度、BFGS 和 L-BFGS 等算法也可以最小化逻辑回归的代价函数。使用这类优化器时，通常需要提供能够返回 $J(\boldsymbol{\theta})$ 与 $\nabla J(\boldsymbol{\theta})$ 的函数。它们常常收敛更快，而且不需要手工指定固定学习率，但内部实现更复杂，实际编程时应优先调用经过验证的数值优化库。
\end{tipbox}

\section{多类别分类}

\begin{definitionbox}[一对多方法]
当 $y\in\{1,2,\ldots,K\}$ 时，可以把一个 $K$ 类问题拆成 $K$ 个二元分类问题。第 $k$ 个分类器把“类别 $k$”作为正类，把其他所有类别作为负类，得到
\[
  h_{\boldsymbol{\theta}^{(k)}}(\mathbf{x})
  =P(y=k\mid\mathbf{x}).
\]
预测时选择输出最大的分类器：
\[
  \hat y
  =\underset{k\in\{1,\ldots,K\}}{\operatorname{arg\,max}}
  \;h_{\boldsymbol{\theta}^{(k)}}(\mathbf{x}).
\]
这种方法称为\concept{一对多（One-vs-Rest）}或一对其余。
\end{definitionbox}

\begin{examplebox}[三类别分类]
若任务需要识别三种类别，则分别训练三个逻辑回归分类器：
\begin{itemize}
  \item 分类器 $1$：类别 $1$ 对类别 $2$、$3$；
  \item 分类器 $2$：类别 $2$ 对类别 $1$、$3$；
  \item 分类器 $3$：类别 $3$ 对类别 $1$、$2$。
\end{itemize}
对新样本同时计算三个输出。例如输出分别为 $0.10$、$0.72$ 和 $0.18$，则预测为类别 $2$。
\end{examplebox}

\section{过拟合与正则化}

\subsection{欠拟合与过拟合}

\begin{definitionbox}[模型拟合的三种状态]
\begin{itemize}
  \item \textbf{欠拟合（Underfitting）：}模型过于简单，连训练数据的主要规律也没有学到，通常表现为高偏差；
  \item \textbf{合适拟合：}模型抓住主要规律，并能较好地推广到未见数据；
  \item \textbf{过拟合（Overfitting）：}模型对训练数据拟合得非常好，甚至代价接近 $0$，却把噪声也当成规律，因而在新数据上表现较差，通常表现为高方差。
\end{itemize}
特征很多、样本较少或多项式次数过高时，过拟合风险尤其明显。
\end{definitionbox}

\begin{keybox}[处理过拟合的基本思路]
常见方法有两类：
\begin{enumerate}
  \item 减少特征数量，人工选择或通过模型选择方法保留真正有用的特征；
  \item 使用正则化，保留特征但限制参数的大小，使模型更平滑、更简单。
\end{enumerate}
删除特征可能损失信息；当多数特征都有一定作用时，正则化通常更合适。
\end{keybox}

\subsection{正则化代价函数}

\begin{definitionbox}[$L_2$ 正则化]
在原代价函数后增加参数平方和作为惩罚项：
\[
  J_{\mathrm{reg}}(\boldsymbol{\theta})
  =J(\boldsymbol{\theta})
  +\frac{\lambda}{2m}\sum_{j=1}^{n}\theta_j^2.
\]
其中 $\lambda\geq0$ 是\concept{正则化参数（Regularization Parameter）}。按照惯例，偏置参数 $\theta_0$ 不参与正则化，因此求和从 $j=1$ 开始。
\end{definitionbox}

\begin{keybox}[正则化参数 $\lambda$ 的作用]
正则化需要在拟合训练数据和限制模型复杂度之间取得平衡：
\begin{itemize}
  \item $\lambda=0$ 时没有正则化，复杂模型可能过拟合；
  \item $\lambda$ 合理时，较大的参数受到抑制，模型通常具有更好的泛化能力；
  \item $\lambda$ 过大时，除偏置外的参数都被压到接近 $0$，模型过于简单并产生欠拟合。
\end{itemize}
$\lambda$ 不应根据测试集表现反复调整，而应使用验证集等模型选择方法确定。
\end{keybox}

\subsection{正则化线性回归}

\begin{definitionbox}[代价函数与梯度下降]
正则化线性回归的代价函数为
\[
  J(\boldsymbol{\theta})
  =\frac{1}{2m}\sum_{i=1}^{m}
  \left(h_{\boldsymbol{\theta}}\!\left(\mathbf{x}^{(i)}\right)-y^{(i)}\right)^2
  +\frac{\lambda}{2m}\sum_{j=1}^{n}\theta_j^2.
\]
偏置和其余参数的更新规则分别为
\begin{align*}
  \theta_0&\leftarrow\theta_0
  -\alpha\frac{1}{m}\sum_{i=1}^{m}
  \left(h_{\boldsymbol{\theta}}\!\left(\mathbf{x}^{(i)}\right)-y^{(i)}\right)x_0^{(i)},\\
  \theta_j&\leftarrow\theta_j
  -\alpha\left[
    \frac{1}{m}\sum_{i=1}^{m}
    \left(h_{\boldsymbol{\theta}}\!\left(\mathbf{x}^{(i)}\right)-y^{(i)}\right)x_j^{(i)}
    +\frac{\lambda}{m}\theta_j
  \right],\quad j\geq1.
\end{align*}
第二个式子也可写成
\[
  \theta_j\leftarrow
  \left(1-\alpha\frac{\lambda}{m}\right)\theta_j
  -\alpha\frac{1}{m}\sum_{i=1}^{m}
  \left(h_{\boldsymbol{\theta}}\!\left(\mathbf{x}^{(i)}\right)-y^{(i)}\right)x_j^{(i)}.
\]
这说明每次更新都会先把参数缩小一点，再沿原代价函数的负梯度方向移动。
\end{definitionbox}

\begin{tipbox}[正则化正规方程]
若使用正规方程，正则化线性回归的解为
\[
  \boldsymbol{\theta}
  =\left(X^TX+\lambda L\right)^{-1}X^T\mathbf{y},
  \qquad
  L=\operatorname{diag}(0,1,\ldots,1).
\]
矩阵 $L$ 左上角为 $0$，保证 $\theta_0$ 不被正则化。实际实现中应使用线性方程求解器或伪逆，而不是显式计算矩阵逆。
\end{tipbox}

\subsection{正则化逻辑回归}

\begin{definitionbox}[正则化逻辑回归的代价函数]
在二元交叉熵后加入同样的平方惩罚项：
\begin{align*}
  J(\boldsymbol{\theta})
  ={}&-\frac{1}{m}\sum_{i=1}^{m}
  \left[
    y^{(i)}\log h_{\boldsymbol{\theta}}\!\left(\mathbf{x}^{(i)}\right)
    +(1-y^{(i)})\log\!\left(1-h_{\boldsymbol{\theta}}\!\left(\mathbf{x}^{(i)}\right)\right)
  \right]\\
  &+\frac{\lambda}{2m}\sum_{j=1}^{n}\theta_j^2.
\end{align*}
其梯度为
\[
  \frac{\partial J}{\partial\theta_0}
  =\frac{1}{m}\sum_{i=1}^{m}
  \left(h_{\boldsymbol{\theta}}\!\left(\mathbf{x}^{(i)}\right)-y^{(i)}\right)x_0^{(i)},
\]
\[
  \frac{\partial J}{\partial\theta_j}
  =\frac{1}{m}\sum_{i=1}^{m}
  \left(h_{\boldsymbol{\theta}}\!\left(\mathbf{x}^{(i)}\right)-y^{(i)}\right)x_j^{(i)}
  +\frac{\lambda}{m}\theta_j,
  \qquad j=1,\ldots,n.
\]
\end{definitionbox}

\begin{mistakebox}[正则化中的常见错误]
\begin{itemize}
  \item 把 $\theta_0$ 加入惩罚项，导致整体预测基线也被不必要地压缩；
  \item 在代价函数中使用 $\lambda/(2m)$，却在梯度中遗漏正则化项的导数 $\lambda\theta_j/m$；
  \item 看到线性回归与逻辑回归的梯度公式外形相似，就误以为二者相同；它们使用的 $h_{\boldsymbol{\theta}}(\mathbf{x})$ 不同；
  \item 认为 $\lambda$ 越大越好。过强的正则化会让模型从过拟合转为欠拟合；
  \item 只比较训练误差来选择模型复杂度和 $\lambda$，从而无法判断模型对新数据的泛化能力。
\end{itemize}
\end{mistakebox}

\section{本周知识脉络}

\begin{keybox}[从概率模型到泛化能力]
第三周的核心流程可以概括为：
\begin{enumerate}
  \item 用 $\boldsymbol{\theta}^{T}\mathbf{x}$ 汇总输入特征；
  \item 用 Sigmoid 函数把结果转成正类概率；
  \item 用阈值和判定边界得到类别预测；
  \item 用二元交叉熵衡量概率预测与真实标签的差异；
  \item 用梯度下降或高级优化算法最小化代价函数；
  \item 对多类别任务训练多个一对多分类器；
  \item 当模型过于复杂时，用正则化限制参数大小，并通过合适的 $\lambda$ 平衡欠拟合与过拟合。
\end{enumerate}
逻辑回归解决的是“如何分类”，正则化进一步解决的是“如何让分类器在未见数据上仍然有效”。
\end{keybox}

\end{document}
